\textbf{Data.} For each seed $s$ we simulate independent trajectories: $32$ training ($1250$ snapshots each), $8$ validation ($500$) and $32$ test ($1500$), all from different random initial conditions, so the splits are disjoint trajectories. The seed determines the data, the network initialisation and the minibatch sequence; the same data seed is shared by all systems, which is why we compare systems with paired tests over seeds. Forecasts start every $50$th snapshot of each test trajectory ($896$ windows per seed) and run for 100 steps ($5$ time units); we report leads $0.5$, 1, 2, 3 and 5 time units as \texttt{t05}\ldots\texttt{t50}, and \texttt{1step} for a single step.

\textbf{Training and tuning.} Every CNN is trained for 1500 Adam steps (batch 32, cosine decay, gradient-norm clip 1). The only tuned hyper-parameter is the peak learning rate, chosen from $\{0.001,0.003\}$ for $K\in\{1,4\}$ in 4 runs on one extra seed (100) by validation RMSE; $0.003$ won for both on the validation RMSE at leads 1 and 2 and was used for every $K$ (no per-$K$ tuning, and no tuning on test). The selected rate is the upper edge of this two-point grid, so a larger rate was not tried and could change the ordering. Persistence and climatology have no parameters; the linear stencil is closed-form with a fixed ridge penalty. The tuning budget is small and $K{=}2,8$ were not tuned separately, so the shared rate may favour $K{=}1,4$; a longer-trained or larger network could also change the picture (Sec.~\ref{sec:limitations}).

\textbf{Registered hypotheses and tests.} Tests were fixed in \texttt{proposal.md} before the runs: paired comparison over 5 seeds (\texttt{rh compare}, paired $t$-test, $\alpha{=}0.05$, no multiplicity correction) of K=4 and K=8 against K=1 for H1 (RMSE at lead 2), stable fraction for H2, variance ratio for H3, ACC lead time of CNN against the baselines for H4. Groups: \texttt{main} (baselines, $K{=}1,4,8$), \texttt{sweep\_K} ($K{=}2$), \texttt{abl\_noise} and \texttt{sweep\_ntrain}.

\textbf{Compute.} Everything ran on one shared CPU machine (macOS, 2 threads per process, at most 2 processes at once), each run within the 6-minute budget. The CNN has \rhval{main/cnn-k-1/l96/n_params/mean:0} parameters. The code is in \texttt{method/}, the driver in \texttt{experiments/run\_all.sh}, and every number in this paper is read from the run registry.
